\begin{afterword}
\summaryhead

A short story. Brennan climbs 256 steps behind a robed guide whose gloves and whisper hide every sign of the guide's sex. At a glass gate the guide asks whether he wants to know; asked what, he answers that it does not matter, the answer is still yes. Inside, a mirrored hall of hooded members of the Bayesian Conspiracy recites the names of four dead probability theorists, each ``dead but not forgotten.''

The guide offers a ring, withdraws it, and asks a question: if three-fourths of the room are women, and three-fourths of the women and half of the men belong to the Heresy of Virtue, what is the probability that the guide, a member, is a man? Brennan says two-elevenths. The guide says one-sixth, the room laughs, and Brennan is urged to give way. He works the answer through aloud, keeps it, and adds that the premises cannot be exactly right, since the room holds an odd number of people. He receives the ring and is made a novice.

\responsehead

The story is fiction and is judged for what it shows. It acts out the proposal of ``To Spread Science, Keep It Secret,'' posted the same day: torches, robes, a secret society, and a test as the price of entry. It adds the initiate's side, which that post only named: shame before the whole room.

What it tests is well chosen. The question is about base rates, and Brennan's answer, two-elevenths, is right. But the test is not the arithmetic. It is whether he keeps a correct answer when an authority overrules him and a crowd laughs, and the guide's invitation to show honesty by giving the answer up is the trap. Brennan earns the ring by keeping a checked answer against an authority and a laughing room, a situation close to Asch's experiments, and by doubting the premises he was given. His doubt is sound. The stated fractions require a room of some multiple of sixteen people, so an odd count cannot fit them. The physical details are right too: complementary plates of glass would cancel each other's distortion, and four mirrored walls do repeat an image in a grid.

What the story shows is the easy case. Brennan can check the answer in his head, so holding firm and being right coincide. The harder question, when to defer to someone who may know more, is not tested here. ``The Ritual,'' earlier in the book, shows a master deferring to a domain expert's counsel. Both are stories, and neither is evidence about which course is right in a given case; this one does not claim to be.

The ending ties the ceremony to the author's essay on Bayes's theorem, placed earlier in the book, which closes: ``You are now an initiate of the Bayesian Conspiracy.'' Here the reader sees what that initiation might look like.

In short: a well-made scene that acts out the previous post's proposal and tests the right things, a base-rate answer and the nerve to keep it under ridicule, with its arithmetic and optics correct; it shows the case where the answer can be checked in one's head, not the harder one where it cannot.
\end{afterword}
